% Source: https://artofproblemsolving.com/community/c3h1945365_2018_amc_8_22
% Stable identity: aops-topic-1945365-q01
\textbf{Similarity.} Since $AB\parallel CD$, the transversal $AC$ makes equal alternate angles $\angle FAB=\angle FCE$, and the transversal $BE$ makes $\angle FBA=\angle FEC$. So $\triangle ABF\sim\triangle CEF$ (AA), and corresponding sides are in the ratio
\[
\dfrac{CE}{AB}=\dfrac{CD}{2\,CD}=\dfrac{1}{2},
\]
because $E$ is the midpoint of $CD$ and $AB=CD$. That is where $1:2$ comes from: it is the ratio of the parallel sides $CE$ and $AB$, which are corresponding sides.

\begin{asy}
usepackage("amsmath");
size(4.5cm);
defaultpen(fontsize(10pt));
pair A=(0,6),B=(6,6),C=(6,0),D=(0,0),E0=(3,0);
pair F=extension(A,C,B,E0);
fill(A--F--E0--D--cycle,gray(0.85));
draw(A--B--C--D--cycle,linewidth(1));draw(A--C);draw(B--E0);
dot(F);dot(E0);
label("$A$",A,NW);label("$B$",B,NE);label("$C$",C,SE);label("$D$",D,SW);label("$E$",E0,S);label("$F$",F,(1.6,0.2));
label("$45$",(1.3,2.2));
label("\mbox{}",(min(currentpicture,user=true)+max(currentpicture,user=true))/2,black); // keeps the text after the figure black when LaTeX typesets the labels inline
\end{asy}

\textbf{Area.} Let the square have side $s$. The heights of $\triangle CEF$ and $\triangle ABF$ to the parallel sides are in ratio $1:2$ and add up to $s$, so the height of $\triangle CEF$ is $\dfrac {s}{3}$ and
\[
[CEF]=\dfrac{1}{2}\cdot\dfrac {s}{2}\cdot\dfrac {s}{3}=\dfrac{s^2}{12}.
\]
Then $[AFED]=[ACD]-[CEF]=\dfrac{s^2}{2}-\dfrac{s^2}{12}=\dfrac{5s^2}{12}=45$, so $s^2=\boxed{\textbf{(B) }108}$.
